\section{Introduction}
Stochastic simulation is a crucial component in the advancement of systems biology modelling. These models look to capture the complex behaviour of intracellular processes with the goal to better understand the functioning of cells as a whole \cite{abkowitzEvidenceThatHematopoiesis1996, arkinStochasticKineticAnalysis1998, kaernStochasticityGeneExpression2005, wilkinsonStochasticModellingQuantitative2009}. It is well defined in the literature that these processes are intrinsically stochastic in nature. Combined with the fact that many living cells often contain low copy numbers of important molecular populations \cite{guptasarmaDoesReplicationinducedTranscription1995, elowitzStochasticGeneExpression2002}, creates high variability and randomness in how these system evolve \cite{kaernStochasticityGeneExpression2005, elowitzStochasticGeneExpression2002, rajNatureNurtureChance2008}. 

To understand how these cellular systems behave, we require complex mathematical models that are able to represent the stochasticity of these biological processes. Stochastic simulation algorithms provide the ability to generate both exact and approximate forward simulations of these models. However, when looking to implement these on realistic biochemical reaction pathways --- involving many chemical species and reactions --- they can quickly become computationally expensive. For biochemical reaction network models, this is due to us requiring the solution to the chemical master equation, the chemical Langevin equation, or reaction rate equations to evolve the system. 

In reality many biochemical systems are analytically intractable, meaning we must rely on computationally expensive approximations to make progress \cite{wilkinsonStochasticModellingQuantitative2009, warneSimulationInferenceAlgorithms2019}. To understand the average behaviour of these high variability reactions, we must generate a large amount of these simulations. Further to this, we often require these results in a timely manner to inform experimental methodology and progress research. As such, high performance computing (HPC) and data centres are leaned upon to yield this data. However, in the advancement of these centres to produce more computational bandwidth, power efficiency has become a secondary priority, often in favour of pure throughput. Recent studies into the greenhouse gas (GHG) emissions caused by computational research has found that the sector produces a non-trivial percentage of global emissions \cite{lannelongueCarbonFootprintEstimation2023, lannelongueGREENERPrinciplesEnvironmentally2023}. Additionally, it is unrealistic to hold the assumption that we can continually add more compute to these centres to reduce wall-time. 

Because of this, we are required to develop more energy-conscious and compute-constrained methods to generate equivalent results. Green high performance computing is a key avenue to explore this problem, by focusing on highly power efficient compute through the use of accelerator hardware \cite{brightwellChallengesHighPerformanceNetworking2010, rostirollaGreenHPCNovelFramework2015}, such as application specific integrated circuits (ASICs), field programmable gate arrays (FPGAs), and tensor processing units (TPUs) \cite{brightwellChallengesHighPerformanceNetworking2010}. In doing so however, we introduce significant complexity in the implementation of these SSAs as this hardware requires more advanced programming techniques to utilise. OpenCL is a heterogeneous computing framework established with the aim to standardise the vendor-specific compilation of these varied devices and increase accessibility to less specialised users \cite{warnePulsecoupledNeuralNetwork2014, warneComparisonHighLevel2014}. 

\subsection{Research Questions}
The focus of this work is to validate and extend mixed-precision multilevel Monte Carlo methods to scalable computing environments. We consider the following research questions: 

\begin{enumerate}
    \item How can mixed-precision coupling techniques for multilevel Monte Carlo be implemented in heterogeneous computing environments?
    \item To what extent can mixed precision multilevel Monte Carlo improve the computational efficiency against standard methods based on double precision arithmetic alone?
\end{enumerate}

\subsection{Contribution}
Through this work we aim to decrease the energy consumption of high performance computing workloads by experimentally validating recent methods proposed to take advantage of mixed precision arithmetic to increase statistical estimator efficiency \cite{haasNestedMLMCFramework2025, kimpsonReducedPrecisionStochasticSimulation2026, bruggerMixedPrecisionMultilevel2014}. We implement Monte Carlo, multilevel Monte Carlo (MLMC), and mixed-precision multilevel Monte Carlo (MP-MLMC) methods using the OpenCL coding language and deploy kernels on a graphics processing unit (GPU). By doing so, we create implementations that can be varied slightly and compiled to a wide range of heterogeneous computing hardware, such as CPUs, GPUs, and FPGAs. With these algorithms we showcase the first externally validated implementation of the methods proposed by Haas and Giles through numerical studies. Furthermore, our OpenCL implementation improves portability and usability of methods in a Systems Biology context. We also introduce an alternative implementation to the proposed coupling scheme using the Box-Muller transformation, substantially reducing the memory footprint of the algorithm, important when considering FPGA implementations.
 
\subsection{Previous Work}
This report forms the second part of the MXB371 and MXB372 research thesis performed over two semesters. The MXB371 report, found in \refsup{sec:supplementary_material}, centred on developing stochastic simulation algorithms (SSAs) using OpenCL. We evaluated these methods on performance and energy consumption for both central processing units (CPUs) and GPUs. This was aimed at understanding varied SSAs and their effects on performance and power utilisation in a heterogeneous environment. Involving the comparison of three methods, two biochemical reaction networks were used to demonstrate results. In this work, we utilise the methods developed for the Euler Maruyama SSA and provide the same two example model definitions. Allowing the focus of this report to be in the utilisation and optimisation of mixed-precision methods for decreasing the computational resources require for accurate simulation estimation. We encourage the reader to view the Literature Review conducted on Biochemical Reaction Networks to gain further insight into the application of our problem. 